% Generated from locale/tr/content/lambda-calculus/lambda-definability/introduction.tex; do not hand-edit.


\olfileid[tr]{lam}{ldf}{int}
\olsection{Giriş}

İlk bakışta lambda hesabı, fonksiyonları ve bunların birbirine uygulanmasını
gösteren ifadelerin son derece soyut bir hesabından ibarettir. Lambda
hesabının sözdiziminde bu fonksiyonların belirli türden nesneler üzerinde
olduğunu düşündüren hiçbir şey yoktur; özellikle sözdizim doğal sayılardan
hiç söz etmez. Temel işlemleri olan uygulama ve lambda soyutlaması, yalnızca
doğal sayılar üzerindeki fonksiyonlara değil, her fonksiyona uygulanır.

Bununla birlikte biraz yaratıcılıkla lambda hesabında aritmetik fonksiyonları,
yani doğal sayılar üzerindeki fonksiyonları tanımlamak mümkündür. Bunun için
her $n\in\Nat$ doğal sayısına, $n$'nin \emph{Church sayısı} olan özel bir
$\lambd$-terimi $\num n$ atarız. Church sayıları adını Alonzo Church'ten alır.

\begin{defn}
  $n\in\Nat$ ise karşılık gelen \emph{Church sayısı} $\num{n}$, $n$'yi
  gösterir:
  \[
    \num{n} \ident \lambd[fx][f^n(x)].
  \]
  Burada $f^n(x)$, $f$'nin $x$'e $n$ kez uygulanmasının sonucudur. Örneğin
  $\num{0}$ terimi $\lambd[fx][x]$, $\num{3}$ terimi ise
  $\lambd[fx][f(f(f\,x))]$'tir.
\end{defn}

Church sayısı $\num n$, iki $f,x$ argümanı alıp $f^n(x)$ döndüren bir
fonksiyonu gösteren lambda terimi olarak kodlanır. Church sayıları açıkça
normal biçimdedir.

Doğal sayıların lambda hesabındaki bir gösterimi, ancak bu sayılarla hesap
yapabiliyorsak yararlıdır. Church sayılarıyla hesap yapmak, bir
$\lambd$-terimi $F$'yi böyle bir $\num n$ Church sayısına uygulayıp birleşik
$F\,\num n$ terimini normal biçime indirgemek demektir. Bu terim her zaman
normal biçime indirgeniyor ve normal biçim her zaman bir $\num m$ Church
sayısı oluyorsa hesaplamanın çıktısını $m$ sayısı olarak düşünebiliriz. Bu
durumda $F$'nin $f\colon\Nat\to\Nat$ fonksiyonunu, yani
$F\,\num n\red\num m$ ancak ve ancak $f(n)=m$ olduğunda geçerli olan
fonksiyonu tanımladığını düşünebiliriz. Church--Rosser özelliği gereği normal
biçimler, varsa, tektir. Dolayısıyla $F\,\num n\red\num m$ ise
$F\,\num n$'nin indirgenebileceği başka bir normal terim, özellikle de başka
bir Church sayısı yoktur.

Tersine, bir $f\colon\Nat\to\Nat$ fonksiyonu verildiğinde $f$'yi bu biçimde
tanımlayan bir $F$ terimi olup olmadığını sorabiliriz. Varsa $F$'nin $f$'yi
\emph{!!{lambda define}s} ve $f$'nin !!{lambda definable} olduğunu söyleriz.
Bunu çok yerli ve kısmi fonksiyonlara genelleyebiliriz.

\begin{defn}
$f\colon\Nat^k\to\Nat$ olsun. Her $n_0$, \dots,~$n_{k-1}$ için
\[
F \, \num{n_0} \, \num{n_1} \dots \num{n_{k-1}} \red
\num{f(n_0,n_1,\dots,n_{k-1})}
\]
$f(n_0,\dots,n_{k-1})$ tanımlı olduğunda geçerli, tanımsız olduğunda ise
$F\,\num{n_0}\,\num{n_1}\dots\num{n_{k-1}}$ normal biçimsizse, $F$ lambda
teriminin $f$'yi \emph{!!{lambda define}s} olduğunu söyleriz.
\end{defn}

Çok basit bir örnek sabit fonksiyonlardır. $C_k\ident\lambd[x][\num{k}]$
terimi, $c_k(n)=k$ olan $c_k\colon\Nat\to\Nat$ fonksiyonunu
!!{lambda define}s. Gerçekten her $n$ için
$C_k\,\num n\ident(\lambd[x][\num{k}])\num n\redone\num{k}$.
Özdeşlik fonksiyonu $\lambd[x][x]$ tarafından !!{lambda defined}.
Daha karmaşık fonksiyonları tanımlamak elbette daha zordur ve çoğu zaman
büyük yaratıcılık gerektirir. Bu nedenle her hesaplanabilir fonksiyonun
!!{lambda definable} olması şaşırtıcı gelebilir. Tersi de doğrudur: bir
fonksiyon !!{lambda definable} ise hesaplanabilirdir.

